\section{Introduction}
\label{sec:intro}

\IEEEPARstart{S}{ecurity} cannot be fully understood by evaluating
autonomous agents independently. An agent may be individually safe,
predictable and robust, yet become a different object once it is placed
in a population of other autonomous agents that observe it, adapt to it,
and are themselves adapting~\cite{park2023generative,wu2023autogen,xi2023rise}.
The resulting feedback loop---attack, observation, defense, counter-adaptation---has
no counterpart in a single-agent red-team trial.

Current work on agentic-AI security has made rapid progress on prompt
injection, tool misuse, memory poisoning, impersonation and
collusion~\cite{greshake2023not,debenedetti2024agentdojo,zhan2024injecagent,chen2024agentpoison,abdelnabi2024cooperation}.
These studies are necessary, but most of them freeze one side of the
interaction. Either the attacker is a script, or the defender is a fixed
policy, or the episode is too short for either side to remember the last
encounter. The question this paper addresses is different:

\begin{quote}
\textit{How does cybersecurity evolve when attackers and defenders are
autonomous learning agents rather than static programs?}
\end{quote}

We treat security as a function of capabilities, goals, interaction,
adaptation and history, rather than as a property of an isolated agent:
\begin{equation}
  \mathrm{Security}
  \;=\;
  f\!\left(\mathcal{C},\, G,\, \mathcal{I},\, \alpha,\, H\right),
  \label{eq:security-fn}
\end{equation}
where $\mathcal{C}$ denotes capabilities, $G$ objectives, $\mathcal{I}$
the interaction structure, $\alpha$ the adaptation regime, and $H$ the
accumulated experience. Equation~\eqref{eq:security-fn} is the conceptual
shift. The benchmark we build is only the instrument.

The contributions are as follows.

\begin{enumerate}
\item \textit{Problem formulation.} We formalise co-evolutionary agent
security as a partially observable, repeated game between learning
populations, and we separate \emph{capability}, \emph{objective} and
\emph{strategy} so that unexpected behaviour can be attributed to
interaction rather than to newly granted tools
(Section~\ref{sec:problem}).
\item \textit{Framework.} We give an agent model, three adaptation
levels (static, episodic, persistent), an interaction graph
$G_t=(V_t,E_t)$, and an operational definition of emergence
(Section~\ref{sec:framework}).
\item \textit{Environment.} We implement a dual-backend cyber
environment---a seedable in-memory simulator and a Kubernetes range---behind
a single interface, with structured tools rather than unconstrained code
execution (Section~\ref{sec:environment}).
\item \textit{Metrics.} In addition to attack success and detection
rates we introduce \emph{co-evolutionary pressure} (CEP) and
\emph{emergent security risk} (ESR), which quantify arms races and
interaction-induced risk (Section~\ref{sec:method}).
\item \textit{Evidence.} A one-versus-one adaptation ladder already
exhibits the phenomenon the framework is built to study: adaptive
attackers evade detection that static attackers cannot, persistent
populations enter an arms race, and several offensive and defensive
behaviours arise without being encoded in objectives
(Section~\ref{sec:results}).
\end{enumerate}

The remainder of the paper is organised as follows.
Section~\ref{sec:related} reviews agentic security, autonomous cyber
defence, multi-agent games and evolutionary computation.
Section~\ref{sec:problem} states the problem.
Section~\ref{sec:framework} presents the framework.
Section~\ref{sec:environment} describes the cyber environment.
Section~\ref{sec:method} details the experimental methodology.
Section~\ref{sec:results} reports the ladder results.
Section~\ref{sec:discussion} discusses implications and limitations.
Section~\ref{sec:conclusion} concludes.
